\chapter{Boolean Analysis — Bernoulli Restriction}
%%% generated by the blueprint pipeline from TCSlib.BooleanAnalysis.Switching.BernoulliRestriction
\section{Overview}

This module sets up the Bernoulli probability model for random restrictions, in which
each of the $n$ variables is independently left free with probability $p$ and fixed to
$\mathsf{true}$ or $\mathsf{false}$ each with probability $(1-p)/2$. It defines the
weight of a single restriction and the probability of an event, and records the basic
facts that the weights are nonnegative, sum to one, and that every event probability is
at most one.

%%%%%%%%%%%%%%%%%%%%%%%%%%%%%%%%%%%%%%%%%%%%%%%%%%%%%%%%%%%%%%%%%%%%%%%%%%%%%%%
\section{Declarations}

\begin{definition}[Bernoulli weight of a restriction]
\lean{SwitchingLemma2.bernoulliRestrWeight}
\statementsource{boolean-ch04-dnf-switching-lmn}{pdf-fc447f326910-p006-b003}
\statementsource{switching-lemma}{pdf-b5e074215b9e-p004-b001}
\label{SwitchingLemma2.bernoulliRestrWeight}
\leanok
\uses{SwitchingLemma2.Restriction, SwitchingLemma2.Restriction.freeVars}
For a parameter $p \in \bbr$ and a restriction $\rho$ on $n$ variables, the Bernoulli
weight is
\[
  p^{\abs{\mathrm{freeVars}(\rho)}} \cdot \left(\tfrac{1-p}{2}\right)^{\,n - \abs{\mathrm{freeVars}(\rho)}},
\]
so each free variable contributes a factor $p$ and each fixed variable a factor
$(1-p)/2$.
\end{definition}

\begin{definition}[Bernoulli event probability]
\lean{SwitchingLemma2.bernoulliRestrProb}
\label{SwitchingLemma2.bernoulliRestrProb}
\leanok
\uses{SwitchingLemma2.Restriction, SwitchingLemma2.bernoulliRestrWeight}
For a parameter $p$ and a (decidable) predicate $\mathrm{event}$ on restrictions, the
probability that the predicate holds under a Bernoulli($p$) random restriction is the
sum over all restrictions $\rho$ of $\mathrm{bernoulliRestrWeight}(p,\rho)$ weighted by
the indicator of $\mathrm{event}(\rho)$.
\end{definition}

\begin{lemma}[Nonnegativity of the Bernoulli weight]
\lean{SwitchingLemma2.bernoulliRestrWeight_nonneg'}
\label{SwitchingLemma2.bernoulliRestrWeight_nonneg'}
\leanok
\uses{SwitchingLemma2.Restriction, SwitchingLemma2.bernoulliRestrWeight}
\difficulty{3}
Let $p \in \bbr$ satisfy $0 \le p \le 1$, and let $\rho$ be a restriction on $n$
variables, with $f$ of its variables left free. Then its Bernoulli weight
$p^{f}\,\bigl(\tfrac{1-p}{2}\bigr)^{\,n - f}$ is nonnegative.
\end{lemma}

\begin{lemma}[Bernoulli restriction weights sum to one]
\lean{SwitchingLemma2.bernoulliRestrWeight_sum_one}
\label{SwitchingLemma2.bernoulliRestrWeight_sum_one}
\leanok
\uses{SwitchingLemma2.Restriction, SwitchingLemma2.Restriction.freeVars, SwitchingLemma2.bernoulliRestrWeight}
\difficulty{5}
Fix $n \in \bbn$ and a real parameter $p$ with $0 \le p \le 1$. For a restriction $\rho$
on $n$ variables, write $F(\rho)$ for its set of free variables. Then, summing the
Bernoulli weight over all restrictions,
\[
  \sum_{\rho} p^{\abs{F(\rho)}} \left(\tfrac{1-p}{2}\right)^{\,n - \abs{F(\rho)}} = 1,
\]
so these weights define a probability distribution on the set of restrictions on $n$
variables.
\end{lemma}

\begin{lemma}[Bernoulli restriction probability is at most one]
\lean{SwitchingLemma2.bernoulliRestrProb_le_one'}
\label{SwitchingLemma2.bernoulliRestrProb_le_one'}
\leanok
\uses{SwitchingLemma2.Restriction, SwitchingLemma2.bernoulliRestrProb, SwitchingLemma2.bernoulliRestrWeight, SwitchingLemma2.bernoulliRestrWeight_nonneg', SwitchingLemma2.bernoulliRestrWeight_sum_one}
\difficulty{4}
Fix $n$ and a real parameter $p$ with $0 \le p \le 1$. Then for every decidable
predicate on restrictions of $n$ variables, the probability that it holds under a
Bernoulli($p$) random restriction is at most $1$.
\end{lemma}
